\section{Exact ordinary moments of continuous resolvent powers}
\label{moment:sec}

The continuous-power transform has a useful specialization that remains
an ordinary convergent integral throughout its stated domain. It gives
every polynomial moment from a regularized Gauss function, and at the
two distinguished imaginary parameters the answer becomes a Gamma
quotient. Generalized harmonic numbers then organize all Bernoulli
moments at once.

For $z\in\mathbb C\setminus\mathbb R$, retain
\[
 \varepsilon=\operatorname{sgn}\operatorname{Im}z,\qquad
 C=-\frac1{z+\varepsilon\pi i},\qquad
 q=\frac{z-\varepsilon\pi i}{z+\varepsilon\pi i},\qquad
 \sigma=-2\varepsilon\pi i.
\]
Thus $|q|<1$. On $0<x<1$ put
\[
 R_z(x)=\frac1{\pi\cot(\pi x)-z},\qquad
 R_z(x)^\lambda=\exp\bigl(\lambda\operatorname{Log}R_z(x)\bigr),
\]
where the principal logarithm is continuous along the interval. Its
argument runs from $0$ to $\varepsilon\pi$. Powers of $C$ use its
principal logarithm as well. In particular,
\[
 R_z(x)^\lambda=C^\lambda
 \left(\frac{1-e^{-\sigma x}}{1-qe^{-\sigma x}}\right)^\lambda,
\]
where the power in parentheses is the radial boundary value of the
holomorphic power that equals one at the origin of the disk.

We use Olver's regularized Gauss function
\begin{equation}\label{moment:eq:regularized}
 \mathbf F(a,b;c;q)=
 \sum_{n=0}^\infty
 \frac{(a)_n(b)_n}{\Gamma(c+n)n!}q^n.
\end{equation}
When $c$ is not a nonpositive integer, this equals
${}_2F_1(a,b;c;q)/\Gamma(c)$. The series definition is essential at
apparent parameter resonances.

\begin{theorem}[Exponential moment generator]
\label{moment:thm:exponential}
For $\operatorname{Re}\lambda>-1$ and every $t\in\mathbb C$, set
$\xi=t/\sigma$. Then
\begin{equation}\label{moment:eq:exponential}
 \boxed{\displaystyle
 \int_0^1 e^{tx}R_z(x)^\lambda\,dx
 =C^\lambda e^{t/2}
   \frac{\Gamma(1+\lambda)}{\Gamma(1+\xi)}
   \mathbf F(\lambda,-\xi;1+\lambda-\xi;q).}
\end{equation}
Both sides are entire functions of $t$. They are holomorphic in
$\lambda$ throughout the half-plane $\operatorname{Re}\lambda>-1$.
\end{theorem}

\begin{proof}
The endpoint expansions are
\[
 R_z(x)=x(1+O(x)),\qquad
 R_z(1-x)=-x(1+O(x)).
\]
The prescribed branches therefore bound the integrand, uniformly on
compact parameter sets, by an integrable constant times
$[x(1-x)]^{\alpha}$, with $\alpha>-1$. Parameter derivatives add
only powers of endpoint logarithms. This proves convergence and the
stated parameter holomorphy for the integral.

Write
\[
 f(w)=\left(\frac{1-w}{1-qw}\right)^\lambda
      =\sum_{n\ge0}d_nw^n,
 \qquad d_0=1.
\]
First suppose $\operatorname{Re}\lambda>0$ and
$\operatorname{Re}\xi<0$. The coefficients are absolutely summable:
the coefficients of $(1-w)^\lambda$ are
$O(n^{-1-\operatorname{Re}\lambda})$, with the terminating cases
understood separately, while $(1-qw)^{-\lambda}$ has geometrically
decaying coefficients. Consequently termwise integration gives
\[
 C^{-\lambda}\int_0^1e^{tx}R_z(x)^\lambda\,dx
 =(e^t-1)\sum_{n\ge0}\frac{d_n}{t-\sigma n}
 =-\frac{e^t-1}{\sigma}\int_0^1
        u^{-\xi-1}(1-u)^\lambda(1-qu)^{-\lambda}\,du.
\]
Euler's integral for the regularized Gauss function
\cite[15.6.1]{DLMF} evaluates the last expression as
\[
 \frac{e^t-1}{t}
 \Gamma(1-\xi)\Gamma(1+\lambda)
 \mathbf F(\lambda,-\xi;1+\lambda-\xi;q).
\]
The Gamma reflection formula and $\sigma=-2\varepsilon\pi i$ give
\[
 \frac{e^t-1}{t}\Gamma(1-\xi)
       =\frac{e^{t/2}}{\Gamma(1+\xi)}.
\]
This proves \eqref{moment:eq:exponential} on the initial open domain.

For $|q|<1$, the series \eqref{moment:eq:regularized} converges
normally on compact subsets of its parameters and is entire in its
three parameters. The right side of \eqref{moment:eq:exponential}
is therefore entire in $t$ and holomorphic for
$\operatorname{Re}\lambda>-1$. The identity theorem first extends
in $t$ and then in $\lambda$, proving the full statement. In
particular, no exclusion of $1+\lambda-\xi=0,-1,-2,\ldots$ is
necessary in the regularized notation.
\end{proof}

The ingredients in this proof are classical Beta and Gauss identities.
The contribution here is their branch-consistent specialization to the
continuous-power resolvent and its exact moment consequences.

\begin{corollary}[Means and logarithmic moments]
\label{moment:cor:mean}
For $\operatorname{Re}\lambda>-1$ and $k\ge0$,
\begin{equation}\label{moment:eq:meanlogs}
 \int_0^1 R_z(x)^\lambda
       \bigl(\operatorname{Log}R_z(x)\bigr)^k\,dx
       =C^\lambda(\operatorname{Log}C)^k.
\end{equation}
In particular, $\int_0^1R_z(x)^\lambda\,dx=C^\lambda$ and
$\int_0^1(\operatorname{Log}R_z(x))^k\,dx=(\operatorname{Log}C)^k$.
\end{corollary}
\begin{proof}
Put $t=0$ in \eqref{moment:eq:exponential}, using
$\mathbf F(\lambda,0;1+\lambda;q)=1/\Gamma(1+\lambda)$.
Differentiate $k$ times in $\lambda$. The endpoint majorants in the
proof of the theorem justify each differentiation under the integral.
\end{proof}

\subsection{The Bernoulli generator and its Gamma specialization}

Let $B_m(x)$ be the Bernoulli polynomials with
$B_1(x)=x-1/2$. Multiplication of
\eqref{moment:eq:exponential} by $t/(e^t-1)$ gives the exact generator
\begin{align}\label{moment:eq:bernoulli-general}
 \sum_{m\ge0}\frac{t^m}{m!}
   \int_0^1 B_m(x)R_z(x)^\lambda\,dx
 ={}&C^\lambda\Gamma(1-t/\sigma)\Gamma(1+\lambda)\notag\\
 &\quad\times
 \mathbf F(\lambda,-t/\sigma;1+\lambda-t/\sigma;q),
 \qquad |t|<2\pi.
\end{align}
The Bernoulli series converges uniformly in $x\in[0,1]$ on every
smaller $t$-disk. The same integrable endpoint majorant justifies the
interchange with the integral. Thus coefficient extraction is an
ordinary-integral identity, including when $\lambda=-1/2$.

\begin{corollary}[A digamma--Lerch first moment]
\label{moment:cor:first}
For $\operatorname{Re}\lambda>-1$,
\begin{align}\label{moment:eq:first}
 \int_0^1\left(x-\frac12\right)R_z(x)^\lambda\,dx
 =\frac{C^\lambda}{\sigma}\bigl[&\psi(1+\lambda)+\gamma
    +q\Phi(q,1,1+\lambda)\notag\\
    &+\operatorname{Log}(1-q)\bigr],
\end{align}
where $\Phi(q,1,a)=\sum_{n\ge0}q^n/(n+a)$.
\end{corollary}
\begin{proof}
With $f(w)$ as above, its nonconstant Fourier coefficients give
\[
 C^{-\lambda}\int_0^1\left(x-\frac12\right)R_z(x)^\lambda dx
       =-\frac1\sigma\int_0^1\frac{f(v)-1}{v}\,dv.
\]
This follows first for $\operatorname{Re}\lambda>0$ by absolute
summation and then for $\operatorname{Re}\lambda>-1$ by holomorphy.
For $0<q<1$, substitute $u=(1-v)/(1-qv)$ in the last integral:
\[
 \int_0^1\frac{f(v)-1}{v}\,dv
 =\int_0^1(u^\lambda-1)
       \left(\frac1{1-u}-\frac q{1-qu}\right)du
 =-\psi(1+\lambda)-\gamma-q\Phi(q,1,1+\lambda)
      -\log(1-q).
\]
The first equality extends to $|q|<1$ by holomorphy. Substitution
proves \eqref{moment:eq:first} with the principal logarithm.
\end{proof}

Define complex generalized harmonic numbers on
$\operatorname{Re}\lambda>-1$ by
\begin{equation}\label{moment:eq:harmonics}
 H_\lambda^{(1)}=\psi(1+\lambda)+\gamma,
 \qquad
 H_\lambda^{(j)}=\zeta(j)-\zeta(j,1+\lambda)\quad(j\ge2).
\end{equation}
For nonnegative integers $\lambda=N$ these are the usual finite
harmonic numbers. Define $h_m(\lambda)$ by
\begin{equation}\label{moment:eq:hgenerator}
 \sum_{m\ge0}h_m(\lambda)u^m
 =\frac{\Gamma(1-u)\Gamma(1+\lambda)}
        {\Gamma(1+\lambda-u)}
 =\exp\left(\sum_{j\ge1}\frac{H_\lambda^{(j)}}j u^j\right)
 \qquad (|u|<\min\{1,|1+\lambda|\}).
\end{equation}
The Taylor series of the Gamma quotient itself is holomorphic for
$|u|<1$. The smaller disk in the exponential representation accounts
for possible zeros of that quotient when $\lambda$ has negative real
part; its logarithmic series must not be continued through such a zero.
Equivalently,
\begin{equation}\label{moment:eq:hrecurrence}
 h_0=1,\qquad
 h_m=\frac1m\sum_{j=1}^m H_\lambda^{(j)}h_{m-j}\quad(m\ge1).
\end{equation}
These are complete homogeneous polynomials in the generalized
harmonic coordinates; for example,
\[
 h_1=H_\lambda^{(1)},\quad
 h_2=\frac{(H_\lambda^{(1)})^2+H_\lambda^{(2)}}2,\quad
 h_3=\frac{(H_\lambda^{(1)})^3+
        3H_\lambda^{(1)}H_\lambda^{(2)}+2H_\lambda^{(3)}}6.
\]

\begin{theorem}[Every Bernoulli moment at the imaginary base points]
\label{moment:thm:bernoulli}
For $z=\varepsilon\pi i$, $\operatorname{Re}\lambda>-1$, and
every integer $m\ge0$,
\begin{equation}\label{moment:eq:allmoments}
 \boxed{\displaystyle
 \int_0^1 B_m(x)
  \left(\frac{\sin(\pi x)}\pi e^{\varepsilon\pi i x}\right)^\lambda dx
 =m!\left(\frac{\varepsilon i}{2\pi}\right)^\lambda
       (-2\varepsilon\pi i)^{-m}h_m(\lambda).}
\end{equation}
The power on the left means
$\exp\{\lambda[\log(\sin(\pi x)/\pi)+\varepsilon\pi i x]\}$.
\end{theorem}
\begin{proof}
At these points $q=0$ and $C=\varepsilon i/(2\pi)$, while
$R_{\varepsilon\pi i}(x)=\sin(\pi x)e^{\varepsilon\pi i x}/\pi$.
Equation \eqref{moment:eq:bernoulli-general} therefore reduces to
$C^\lambda$ times the Gamma quotient
\eqref{moment:eq:hgenerator} at $u=t/\sigma$. Comparing coefficients
proves the formula. The logarithmic Taylor expansion of that quotient
gives \eqref{moment:eq:hgenerator} from
\eqref{moment:eq:harmonics}; logarithmic differentiation then gives
\eqref{moment:eq:hrecurrence}.
\end{proof}

At the negative half power,
\[
 H_{-1/2}^{(1)}=-2\log2,\qquad
 H_{-1/2}^{(j)}=-(2^j-2)\zeta(j)\quad(j\ge2).
\]
The first two nonconstant moments are consequently
\begin{align}\label{moment:eq:halfexamples}
 \int_0^1\left(x-\frac12\right)
  \sqrt{\frac\pi{\sin(\pi x)}}e^{-\varepsilon\pi i x/2}\,dx
 &=-\frac{1+\varepsilon i}{\sqrt\pi}\log2,\notag\\
 \int_0^1\left(x^2-x+\frac16\right)
  \sqrt{\frac\pi{\sin(\pi x)}}e^{-\varepsilon\pi i x/2}\,dx
 &=\sqrt\pi(1-\varepsilon i)
     \left(\frac1{12}-\frac{\log^22}{\pi^2}\right).
\end{align}
Both endpoints have only inverse-square-root singularities, so these
are absolutely convergent ordinary integrals. For example the second
identity also yields the entirely real evaluation
\begin{equation}\label{moment:eq:realhalf}
 \int_0^1\left(x^2-x+\frac16\right)
 \frac{\cos(\pi x/2)}{\sqrt{\sin(\pi x)}}\,dx
 =\frac1{12}-\frac{\log^22}{\pi^2}.
\end{equation}
At the positive half power,
$H_{1/2}^{(1)}=2-2\log2$ and
$H_{1/2}^{(j)}=2^j-(2^j-2)\zeta(j)$ for $j\ge2$.
Thus every moment at either half power belongs to an explicitly
specified algebra of zeta values, $\log2$, $\pi$, and the displayed
square-root prefactor. Further $\lambda$-derivatives of
\eqref{moment:eq:allmoments} give all corresponding moments with
arbitrary powers of $\operatorname{Log}R_{\varepsilon\pi i}$.

\paragraph{Verification.}
The independent script \texttt{check\_moments.py} compares the
Gamma--Gauss generator against direct real-interval quadrature in both
half-planes, including nonintegral powers with negative real part and
powers with real part greater than one. It also compares four Bernoulli
moments with \eqref{moment:eq:allmoments}, reconstructs the first nine
Fourier coefficients by two exact symbolic formulas, and rejects a
changed coefficient. Its separate spectral and finite-part checks audit
the conventions used in the surrounding continuous-power transform.
These numerical comparisons are diagnostics, not interval certificates
or steps in the proofs above.
